
\documentclass[12pt,a4paper]{article}

\usepackage[utf8]{inputenc}
\usepackage[T1]{fontenc}
\usepackage[brazil]{babel}
\usepackage{lmodern}
\usepackage[a4paper,top=3cm,left=3cm,right=2cm,bottom=2cm]{geometry}
\usepackage{setspace}
\usepackage{xcolor}
\usepackage{graphicx}
\usepackage{booktabs}
\usepackage{tabularx}
\usepackage{listings}
\usepackage{hyperref}
\usepackage{float}

\onehalfspacing
\hypersetup{
    colorlinks=true,
    linkcolor=black,
    urlcolor=blue
}

\lstset{
    basicstyle=\ttfamily\footnotesize,
    backgroundcolor=\color{gray!8},
    frame=single,
    rulecolor=\color{gray!40},
    breaklines=true,
    columns=fullflexible,
    keepspaces=true,
    showstringspaces=false,
    literate=
        {á}{{\'a}}1 {é}{{\'e}}1 {í}{{\'i}}1
        {ó}{{\'o}}1 {ú}{{\'u}}1 {â}{{\^a}}1
        {ê}{{\^e}}1 {ô}{{\^o}}1 {ã}{{\~a}}1
        {õ}{{\~o}}1 {ç}{{\c{c}}}1
        {Á}{{\'A}}1 {É}{{\'E}}1
        {à}{{\`a}}1
}

\begin{document}

\begin{titlepage}
    \centering

    {\large
    INSTITUTO FEDERAL DE EDUCAÇÃO, CIÊNCIA E TECNOLOGIA DO CEARÁ\\
    CAMPUS BOA VIAGEM\\
    TECNOLOGIA EM ANÁLISE E DESENVOLVIMENTO DE SISTEMAS
    \par}

    \vspace{3cm}

    {\large
    FERNANDO EVERSON FERREIRA CAVALCANTE
    \par}

    \vspace{3cm}

    {\Large\bfseries
    TRANSPORTE ESCOLAR
    \par}

    \vspace{0.5cm}

    {\large
    Sistema web orientado a objetos em PHP
    \par}

    \vfill

    {\large
    Relatório do trabalho individual da disciplina Programação Web I\\
    Professor: Prof. Dr. Renato William Rodrigues de Souza
    \par}

    \vspace{2cm}

    {\large
    Boa Viagem, CE\\
    2026
    \par}
\end{titlepage}

\tableofcontents
\newpage

\section{Introdução}

O transporte escolar é um serviço que necessita de organização para relacionar veículos, rotas e alunos. O controle dessas informações permite identificar os estudantes atendidos por determinada rota, verificar a capacidade dos veículos e acompanhar a quantidade de vagas ocupadas.

O sistema \textbf{Transporte Escolar} foi desenvolvido com o objetivo de representar essas operações em uma aplicação web simples, utilizando PHP e conceitos de Programação Orientada a Objetos (POO).

A aplicação permite cadastrar ônibus e vans, criar rotas associadas aos veículos disponíveis e cadastrar alunos nas rotas existentes. Também oferece páginas para consultar os registros e visualizar a ocupação de cada veículo.

O público-alvo do sistema são pessoas responsáveis pela organização do transporte escolar, como funcionários administrativos e gestores que necessitam consultar informações sobre veículos e estudantes.

O desenvolvimento foi realizado utilizando PHP, HTML, CSS e JSON, sem a utilização de frameworks ou sistemas gerenciadores de banco de dados. Os registros são armazenados em um arquivo JSON, permitindo sua recuperação após o encerramento do servidor.

Este relatório apresenta os requisitos do sistema, a modelagem das classes, a implementação dos conceitos de orientação a objetos, a camada web, os testes, as ferramentas utilizadas, o uso de inteligência artificial e as dificuldades encontradas durante o desenvolvimento.

\section{Requisitos}

\subsection{Funcionalidades}

O sistema possui as seguintes funcionalidades:

\begin{itemize}
    \item Exibir uma página inicial com acesso às funcionalidades.
    \item Cadastrar ônibus e vans.
    \item Listar veículos cadastrados.
    \item Cadastrar rotas associadas a veículos existentes.
    \item Listar rotas e consultar seus detalhes.
    \item Cadastrar alunos em rotas selecionadas.
    \item Listar os alunos cadastrados e suas respectivas rotas.
    \item Calcular a ocupação dos veículos.
    \item Validar os dados recebidos pelos formulários.
    \item Armazenar e recuperar informações utilizando JSON.
    \item Disponibilizar os registros por meio de uma resposta HTTP em formato JSON.
    \item Retornar HTTP 404 para páginas ou identificadores de rotas inexistentes.
\end{itemize}

A navegação é controlada pelo parâmetro \texttt{pagina} da URL. Os formulários utilizam o método HTTP POST, enquanto as consultas são realizadas por GET.

\subsection{Regras de negócio}

As principais regras implementadas diretamente nas classes são apresentadas na Tabela~\ref{tab:regras}.

\begin{table}[H]
\centering
\small
\begin{tabularx}{\linewidth}{@{}cXX@{}}
\toprule
\textbf{Nº} & \textbf{Regra} & \textbf{Classe e método} \\
\midrule
1 & A capacidade do veículo deve ser maior que zero. &
\texttt{Veiculo::\_\_construct()} \\
2 & Um aluno não pode possuir matrícula duplicada na mesma rota. &
\texttt{Rota::adicionarAluno()} \\
3 & A quantidade de alunos não pode ultrapassar a capacidade do veículo. &
\texttt{Rota::adicionarAluno()} \\
\bottomrule
\end{tabularx}
\caption{Regras de negócio implementadas nas classes.}
\label{tab:regras}
\end{table}

Além dessas regras, a camada web realiza outras validações:

\begin{itemize}
    \item Impede o cadastro de veículos com placas repetidas.
    \item Impede o cadastro de rotas com nomes repetidos.
    \item Impede a associação de um veículo a duas rotas.
    \item Impede matrículas duplicadas entre as rotas.
    \item Exige a seleção de uma rota existente no cadastro de alunos.
    \item Verifica o preenchimento dos campos obrigatórios.
\end{itemize}

Quando uma regra é violada, o sistema utiliza exceções para interromper a operação e apresentar uma mensagem de erro.

\section{Modelagem}

A modelagem foi organizada em cinco classes principais: \texttt{Veiculo}, \texttt{Onibus}, \texttt{Van}, \texttt{Aluno} e \texttt{Rota}.

A classe \texttt{Veiculo} representa as características compartilhadas pelos veículos. As classes \texttt{Onibus} e \texttt{Van} especializam esse comportamento. A classe \texttt{Aluno} representa os estudantes, enquanto \texttt{Rota} organiza o veículo utilizado e os alunos associados.

\begin{figure}[H]
    \centering
    \IfFileExists{diagrama.png}{
        \includegraphics[width=0.95\linewidth]{diagrama.png}
    }{
        \fbox{
            \parbox{0.85\linewidth}{
                \centering
                \vspace{1.5cm}
                Diagrama de classes\\
                Adicionar o arquivo diagrama.png ao Overleaf.
                \vspace{1.5cm}
            }
        }
    }
    \caption{Diagrama de classes do sistema Transporte Escolar.}
    \label{fig:diagrama}
\end{figure}

\subsection{Herança}

A herança é aplicada por meio da classe mãe \texttt{Veiculo} e das classes filhas \texttt{Onibus} e \texttt{Van}.

Essa relação segue o princípio ``é um'', pois tanto um ônibus quanto uma van são tipos de veículos.

A classe mãe concentra atributos e métodos comuns, como placa, modelo, capacidade e cálculo de lotação.

As classes filhas utilizam \texttt{extends} e chamam o construtor da classe mãe por meio de \texttt{parent::\_\_construct()}.

O método \texttt{calcularLotacao()} é sobrescrito para apresentar o tipo de veículo no resultado.

\subsection{Composição}

A classe \texttt{Rota} possui um objeto \texttt{Veiculo} e mantém uma lista de objetos \texttt{Aluno}.

Essa organização representa a relação ``tem um'', pois uma rota possui um veículo associado e reúne os alunos atendidos.

O método \texttt{adicionarAluno()} controla a inclusão dos estudantes e verifica as regras de matrícula duplicada e capacidade.

O método \texttt{calcularOcupacao()} utiliza o veículo associado para apresentar a quantidade de vagas ocupadas.

Os veículos também podem existir independentemente das rotas, pois são cadastrados separadamente. Portanto, a associação entre rota e veículo não exige que o veículo seja criado dentro da classe \texttt{Rota}.

\section{Implementação}

\subsection{Encapsulamento e construtores (OO1 e OO4)}

O encapsulamento foi aplicado por meio de atributos privados nas classes.

O exemplo abaixo apresenta parte da implementação de \texttt{Veiculo}.

\begin{lstlisting}[language=PHP]
class Veiculo
{
    private string $placa;
    private string $modelo;
    private int $capacidade;

    public function __construct(
        string $placa,
        string $modelo,
        int $capacidade
    ) {
        if ($capacidade < 1) {
            throw new Exception(
                "A capacidade do veículo deve ser maior que zero."
            );
        }

        $this->placa = $placa;
        $this->modelo = $modelo;
        $this->capacidade = $capacidade;
    }
}
\end{lstlisting}

O uso de \texttt{private} impede o acesso direto aos atributos fora da classe. Os valores são recebidos pelo construtor e podem ser consultados por meio dos getters.

O construtor também realiza uma validação para impedir a criação de veículos com capacidade inválida.

Todas as cinco classes possuem construtores e atributos tipados, e os arquivos utilizam \texttt{declare(strict\_types=1)}.

\subsection{Herança e sobrescrita (OO2)}

A classe \texttt{Onibus} herda de \texttt{Veiculo} e sobrescreve o método de cálculo de lotação.

\begin{lstlisting}[language=PHP]
class Onibus extends Veiculo
{
    public function __construct(
        string $placa,
        string $modelo,
        int $capacidade
    ) {
        parent::__construct(
            $placa,
            $modelo,
            $capacidade
        );
    }

    public function calcularLotacao(
        int $quantidadeAlunos
    ): string {
        return "Ônibus: " .
            parent::calcularLotacao($quantidadeAlunos);
    }
}
\end{lstlisting}

A classe \texttt{Van} utiliza a mesma estrutura de herança, apresentando a identificação ``Van'' no resultado.

A sobrescrita permite especializar o comportamento sem repetir toda a lógica de cálculo da classe mãe.

\subsection{Polimorfismo (OO3)}

O polimorfismo é demonstrado no arquivo \texttt{teste.php}, por meio de um array contendo objetos de diferentes classes.

\begin{lstlisting}[language=PHP]
$veiculos = [$onibus, $van];

foreach ($veiculos as $veiculo) {
    echo $veiculo->calcularLotacao(1) . "\n";
}
\end{lstlisting}

Mesmo utilizando a mesma chamada de método, o resultado depende do tipo concreto do objeto.

\begin{lstlisting}
Ônibus: 1/2 vagas ocupadas
Van: 1/8 vagas ocupadas
\end{lstlisting}

Esse comportamento permite trabalhar com diferentes tipos de veículos por meio de uma interface comum herdada de \texttt{Veiculo}.

\subsection{Regras de negócio com exceções (OO5)}

O método \texttt{adicionarAluno()} verifica se a matrícula já existe na rota e se ainda há vagas disponíveis.

\begin{lstlisting}[language=PHP]
public function adicionarAluno(Aluno $aluno): void
{
    foreach ($this->alunos as $alunoAdicionado) {
        if (
            $alunoAdicionado->getMatricula() ===
            $aluno->getMatricula()
        ) {
            throw new Exception(
                "Este aluno já está cadastrado na rota."
            );
        }
    }

    if (
        count($this->alunos) >=
        $this->veiculo->getCapacidade()
    ) {
        throw new Exception(
            "A capacidade máxima do veículo foi atingida."
        );
    }

    $this->alunos[] = $aluno;
}
\end{lstlisting}

O uso de \texttt{throw new Exception()} permite interromper operações que violam as regras do sistema.

Na camada web, essas exceções são tratadas com \texttt{try/catch}, permitindo apresentar uma mensagem ao usuário.

\subsection{Composição (OO6)}

A classe \texttt{Rota} mantém referências para objetos de outras classes.

\begin{lstlisting}[language=PHP]
class Rota
{
    private string $nome;
    private Veiculo $veiculo;
    private array $alunos = [];

    public function __construct(
        string $nome,
        Veiculo $veiculo
    ) {
        $this->nome = $nome;
        $this->veiculo = $veiculo;
    }

    public function calcularOcupacao(): string
    {
        return $this->veiculo->calcularLotacao(
            count($this->alunos)
        );
    }
}
\end{lstlisting}

A rota utiliza a capacidade do veículo associado e a quantidade de alunos armazenados para calcular a ocupação.

Dessa forma, a operação depende da colaboração entre os objetos.

\subsection{Organização em arquivos e conversão para JSON (OO7 e OO8)}

As classes estão separadas na pasta \texttt{src/}.

\begin{lstlisting}
src/
    Aluno.php
    Veiculo.php
    Onibus.php
    Van.php
    Rota.php
\end{lstlisting}

O arquivo \texttt{public/index.php} carrega as classes utilizando \texttt{require\_once}.

\begin{lstlisting}[language=PHP]
require_once __DIR__ . '/../src/Veiculo.php';
require_once __DIR__ . '/../src/Onibus.php';
require_once __DIR__ . '/../src/Van.php';
require_once __DIR__ . '/../src/Aluno.php';
require_once __DIR__ . '/../src/Rota.php';
\end{lstlisting}

Todas as classes possuem o método \texttt{paraArray()}, que permite converter os objetos para arrays associativos.

\begin{lstlisting}[language=PHP]
public function paraArray(): array
{
    return [
        "placa" => $this->placa,
        "modelo" => $this->modelo,
        "capacidade" => $this->capacidade
    ];
}
\end{lstlisting}

Os arrays são utilizados na geração das respostas JSON e no armazenamento dos dados.

A persistência utiliza \texttt{json\_encode()}, \texttt{json\_decode()}, \texttt{file\_put\_contents()} e \texttt{file\_get\_contents()}.

O arquivo \texttt{dados/sistema.json} contém os veículos e as rotas, incluindo os alunos associados.

\section{Camada web}

A camada web foi implementada no arquivo \texttt{public/index.php}, responsável pelo processamento das requisições e pela apresentação das páginas.

\subsection{Navegação}

A navegação utiliza o parâmetro \texttt{pagina} da URL.

As principais páginas são:

\begin{itemize}
    \item \texttt{?pagina=inicio}
    \item \texttt{?pagina=veiculos}
    \item \texttt{?pagina=cadastro-veiculo}
    \item \texttt{?pagina=rotas}
    \item \texttt{?pagina=cadastro-rota}
    \item \texttt{?pagina=alunos}
    \item \texttt{?pagina=cadastro}
\end{itemize}

A consulta aos detalhes de uma rota utiliza também o parâmetro \texttt{id}.

\subsection{Tratamento de páginas inexistentes}

O sistema mantém uma lista de páginas permitidas.

Caso o usuário solicite uma página não reconhecida, é utilizado:

\begin{lstlisting}[language=PHP]
if (!in_array($pagina, $paginasPermitidas, true)) {
    http_response_code(404);
    $pagina = '404';
}
\end{lstlisting}

A mesma resposta é utilizada quando o identificador de uma rota não existe.

\subsection{Formulários POST}

Os formulários permitem cadastrar veículos, rotas e alunos.

O processamento acontece no início do arquivo PHP, antes da apresentação do HTML.

As informações são obtidas por meio de \texttt{\$\_POST}, validadas e utilizadas na criação dos objetos.

Um exemplo de validação é a verificação de matrícula duplicada.

\begin{lstlisting}[language=PHP]
foreach ($rotas as $rotaAtual) {
    foreach ($rotaAtual->getAlunos() as $alunoAtual) {
        if ($alunoAtual->getMatricula() === $matricula) {
            throw new Exception(
                'Este aluno já está cadastrado no sistema.'
            );
        }
    }
}
\end{lstlisting}

As mensagens de erro são apresentadas na interface após o tratamento da exceção.

\subsection{Resposta JSON}

O sistema permite consultar os dados utilizando o parâmetro \texttt{formato=json}.

A resposta utiliza o cabeçalho:

\begin{lstlisting}[language=PHP]
header('Content-Type: application/json; charset=utf-8');
\end{lstlisting}

O corpo contém as coleções \texttt{veiculos} e \texttt{rotas}.

Um exemplo simplificado da estrutura é:

\begin{lstlisting}
{
    "veiculos": {
        "ABC-1234": {
            "placa": "ABC-1234",
            "modelo": "Mercedes-Benz",
            "capacidade": 20,
            "tipo": "Ônibus"
        }
    },
    "rotas": {
        "centro": {
            "nome": "Rota Centro",
            "veiculo": {},
            "alunos": [],
            "ocupacao": "Ônibus: 0/20 vagas ocupadas"
        }
    }
}
\end{lstlisting}

O exemplo acima representa a estrutura da resposta, não uma captura literal da execução.

\subsection{Proteção da saída HTML}

Os dados exibidos nas páginas são tratados com \texttt{htmlspecialchars()}.

\begin{lstlisting}[language=PHP]
$esc = static function (string $valor): string {
    return htmlspecialchars(
        $valor,
        ENT_QUOTES,
        'UTF-8'
    );
};
\end{lstlisting}

Esse tratamento evita que caracteres especiais presentes nos dados sejam interpretados como marcação HTML.

\section{Testes}

Os testes foram realizados por meio do arquivo \texttt{teste.php}, de requisições HTTP e de verificações manuais no navegador.

\subsection{Testes das classes}

O comando utilizado é:

\begin{lstlisting}
php teste.php
\end{lstlisting}

A saída esperada para os testes principais é:

\begin{lstlisting}
=== POLIMORFISMO ===
Ônibus: 1/2 vagas ocupadas
Van: 1/8 vagas ocupadas

=== CADASTRO DE ALUNOS ===
Ana adicionada com sucesso.
Carlos adicionado com sucesso.

=== TESTE DE ALUNO DUPLICADO ===
Erro esperado: Este aluno já está cadastrado na rota.

=== TESTE DE CAPACIDADE ===
Erro esperado: A capacidade máxima do veículo foi atingida.

=== TESTE DE VEÍCULO INVÁLIDO ===
Erro esperado: A capacidade do veículo deve ser maior que zero.
\end{lstlisting}

O primeiro teste demonstra polimorfismo, pois diferentes objetos executam o mesmo método.

Os testes seguintes verificam o cadastro de alunos e a aplicação das regras de negócio.

A recusa de matrícula duplicada demonstra o funcionamento da validação implementada em \texttt{Rota}.

O teste de capacidade comprova que a rota impede a inclusão de alunos acima do limite do veículo.

Por fim, a criação de um veículo com capacidade zero é recusada pelo construtor.

\subsection{Testes HTTP}

As requisições podem ser executadas utilizando o terminal e o programa \texttt{curl}.

\begin{lstlisting}
curl -i "http://localhost:8000/?pagina=rotas"
\end{lstlisting}

Resultado esperado: HTTP 200.

\begin{lstlisting}
curl -i "http://localhost:8000/?pagina=inexistente"
\end{lstlisting}

Resultado esperado: HTTP 404.

\begin{lstlisting}
curl -i "http://localhost:8000/?formato=json"
\end{lstlisting}

Resultado esperado: HTTP 200, com o cabeçalho:

\begin{lstlisting}
Content-Type: application/json; charset=utf-8
\end{lstlisting}

Os comandos acima também estão documentados em \texttt{docs/testes.md}.

\subsection{Testes dos cadastros}

Após a ampliação do sistema, foram realizadas verificações manuais dos cadastros de veículos, rotas e alunos.

O cadastro de veículos permite selecionar ônibus ou van, informar placa, modelo e capacidade.

O cadastro de rotas permite selecionar um veículo existente.

O cadastro de alunos permite escolher a rota à qual o estudante será associado.

As funcionalidades foram verificadas no navegador durante o desenvolvimento.

\subsection{Persistência}

O armazenamento em JSON permite recuperar os registros após novas requisições e reinicialização do servidor.

O arquivo é atualizado quando um cadastro é realizado com sucesso.

Essa solução foi escolhida pela simplicidade e por atender ao escopo acadêmico do projeto, embora não ofereça os recursos de consulta e controle de concorrência de um banco de dados.

\section{Ferramentas utilizadas}

\begin{table}[H]
\centering
\small
\begin{tabularx}{\linewidth}{@{}lX@{}}
\toprule
\textbf{Ferramenta} & \textbf{Para que foi usada} \\
\midrule
PHP 8.5 & Desenvolvimento das classes, regras de negócio e servidor local. \\
VS Code & Escrita e edição dos arquivos do projeto. \\
Git e GitHub & Controle de versões e armazenamento do código-fonte. \\
curl & Verificação das respostas HTTP. \\
Navegador e DevTools & Visualização e testes da interface web. \\
ChatGPT & Apoio na implementação, resolução de erros e documentação. \\
Gravador de tela & Ferramenta a informar após a gravação do vídeo. \\
\bottomrule
\end{tabularx}
\caption{Ferramentas utilizadas no desenvolvimento.}
\end{table}

\section{Uso de LLM}

Durante o desenvolvimento do projeto, o ChatGPT foi utilizado como ferramenta de apoio para compreender conceitos de orientação a objetos, organizar a estrutura das classes, implementar funcionalidades e solucionar problemas.

A ferramenta também auxiliou na implementação da persistência em JSON, na ampliação dos cadastros e na revisão da documentação.

Uma das principais contribuições foi a possibilidade de analisar mensagens de erro e compreender suas causas.

Um exemplo ocorreu quando a execução de \texttt{php -l} apresentou a mensagem de que \texttt{strict\_types} precisava ser a primeira instrução do arquivo. A partir da explicação recebida, a organização inicial do código foi corrigida.

Entretanto, a utilização da LLM também apresentou limitações. Durante a elaboração da documentação, uma mensagem da saída do terminal foi reproduzida incorretamente, demonstrando que respostas geradas automaticamente podem conter diferenças em relação aos resultados reais.

Além disso, códigos sugeridos precisam ser testados antes de serem considerados corretos. Uma implementação aparentemente adequada pode apresentar problemas de sintaxe, validação ou comportamento durante a execução.

Dessa forma, a LLM foi utilizada como apoio, e não como substituta dos testes e da análise do código.

O histórico das interações está documentado no arquivo \texttt{docs/llm.md}.

\section{Dificuldades, aprendizados e conclusão}

Uma das dificuldades encontradas foi organizar a aplicação utilizando conceitos de orientação a objetos, especialmente ao definir quais responsabilidades deveriam pertencer a cada classe.

A separação entre \texttt{Veiculo}, \texttt{Onibus}, \texttt{Van}, \texttt{Aluno} e \texttt{Rota} ajudou a compreender melhor a utilização de herança, encapsulamento e colaboração entre objetos.

Outra dificuldade foi integrar as classes à camada web. Foi necessário compreender o recebimento de informações pelos formulários, a validação dos dados e o tratamento das exceções.

A persistência também exigiu atenção, pois os objetos PHP precisavam ser convertidos em arrays e posteriormente reconstruídos a partir do arquivo JSON.

Com a ampliação do sistema, tornou-se necessário adaptar a estrutura de armazenamento para permitir o cadastro independente de veículos e sua associação às rotas.

Durante esse processo, foram reforçados conhecimentos sobre construtores, métodos, atributos privados, herança, polimorfismo, tratamento de exceções, requisições HTTP e manipulação de arquivos.

O sistema desenvolvido atende ao objetivo de organizar informações básicas do transporte escolar, permitindo cadastrar veículos, rotas e alunos, consultar os registros e verificar a ocupação dos veículos.

Como melhorias futuras, poderiam ser implementadas funcionalidades de edição e exclusão, autenticação de usuários e armazenamento em banco de dados.

Conclui-se que o projeto contribuiu para a aplicação prática dos conceitos estudados na disciplina de Programação Web I, demonstrando como a orientação a objetos pode ser utilizada no desenvolvimento de uma aplicação web funcional.

\section*{Links}

\begin{itemize}
    \item Repositório:
    \url{https://github.com/fernando-ferreira-caval/pweb1-trabalho-oo}

    \item Vídeo de apresentação:
    \textit{Adicionar link após a gravação.}
\end{itemize}

\end{document}
